\section{Discussion}
\label{sec:discussion}

\subsection{What the evidence establishes}
\label{sec:discussion-establishes}

The evaluation supports four claims at graded strength. Device-level
integration is established strongly (RQ1): five protocol stacks behave
identically on platform and device sides, including under faults, and the
behavior survives configuration changes and re-runs. Mechanism
attribution is established at the platform tier (RQ2). Its
epistemic status should be stated plainly: the ablation is an engineering
validation that the design's attribution story is true (the interlock
predicate does carry interception; the structural layers survive its
removal), not a causal comparison of governed against ungoverned
operation on plant outcomes. Such a comparison was deliberately not run
on a live path, and plant-KPI superiority is not claimed. What the
ablation does establish is defense in depth (the hard range survives
every disabled configuration) and the cost of governance on this
workload. The cost result is zero false blocks on injected legitimate
traffic. It is not a false-block rate for natural or adversarially
ambiguous plant traffic, and not a statement about in-range harmful
writes, such as compliant amplitudes applied at harmful times, which the
interval layers do not address. The closed-loop result (RQ3) is a feasibility demonstration
with unusually complete instrumentation: every decision is
ledger-reconstructed, approvals bind concrete tuples, and an independent
physical recheck corroborates the outcome. It is not a statistical
comparison, and one of its two episodes exercised the human decision
point with an automated approver. The adversarial result (RQ5) is a
negative-outcome containment case: it shows the platform bounding a weak
model's damage to zero writes, not that weak models can be made to
succeed. Throughout, the model-driven tier's benefits (RQ4) are
simulation-tier results by construction.

\subsection{Implications for practitioners}
\label{sec:discussion-practitioners}

Three lessons generalize beyond this implementation. First,
\emph{governance must live in configuration, not in prompts}. The same
agent surface that passed the ablation's legitimate traffic and blocked
its attacks never received a rule as text; the rules were intersections
of configuration facts. Teams deploying LLM agents against plants should
resist encoding limits in system prompts, which are advisory, and
instead make them admission predicates, which are enforceable.

Second, \emph{attribution is a first-class requirement, not a logging
by-product}. The ablation could identify the load-bearing layer only
because every layer records its own verdict. Operators debugging a
blocked write, or an auditor reconstructing an executed one, need
the excluding layer and the configuration revision, and they need
rejections recorded as faithfully as executions. The
\emph{attribution $\neq$ prevention} separation also has a practical
consequence: disabling a soft layer for commissioning does not blind the
audit, and a write that executes under a disabled interlock is still
visible.

Third, \emph{receipt grading changes what agents can honestly report}.
An agent that cites a readback-verified write and one that cites a
transport acknowledgment are making different claims; surfacing the
grade in the tool result makes the difference visible to the model and
to any human reading the report. This is the same discipline
human-factors research asks of automation: the human decision should be
exercised against concrete, inspectable evidence
\cite{parasuraman2000model}, not against a summary. The
complacency failure mode \cite{parasuramanmanzey2010complacency}
argues for keeping the approval payload specific enough that approving
is an informed act.

A deployment path follows from the architecture: commission bindings in
manual mode with batch-window interlocks on, review the ledger's
rejection records to tune limits, and promote individual bindings to
auto only after their rejection history is understood. The
configuration that the ablation exercised maps directly onto this
ramp-up.

For an application-oriented venue, the relevant question is whether the
governance layer survives deployment, not whether the underlying model is
stronger. The evaluation is deliberately staged: protocol-level fidelity
on real stacks, process-level behavior on a simulator, and plant-level
generalization as the next empirical step. The staging is a scope choice,
not a claim that simulation substitutes for a plant. A plant pilot should
start read-only, then move to single-point manual writes, then to bounded
automatic writes, keeping the same ledger and approval records throughout;
the protocol matrix and the four-layer path are the parts that transfer
unchanged.

\subsection{Limitations and threats to validity}
\label{sec:discussion-limits}

\textbf{Construct validity.} The ablation measures interception of
injected policy-breaching writes and false blocks on legitimate ones;
these are mechanistic constructs, not plant-outcome constructs. We do
not claim that governed supervision improves plant KPIs relative to
ungoverned supervision: that comparison was deliberately not run on a
live path, and the model-tier benefits are simulation-tier.

\textbf{External validity.} All plant contact is against simulator
backends: a physical-engine line simulator and a PLC node simulator
exposing real protocol stacks. The protocol behavior is real; the
process physics is simulated. The case-study line is single-site, and
the LLM episodes use one 9B-class open model family; stronger or
weaker models may shift where governance is exercised, though the
containment episode suggests the bound holds when the model is weak.

\textbf{Internal validity.} Case-study timelines are reconstructed from
the platform's own ledger; an independent log source does not exist for
these episodes (the physical recheck provides one external corroboration
point). The ablation's workload is synthetic by design, which is what
makes its attribution exact. Statistical-conclusion validity is
deliberately thin: most reported numbers are single-run or small-$n$
descriptive statistics, and we have marked each as such rather than
attaching tests they cannot support. Experimenter bias is a live
threat: the attack sets and legitimate workloads were designed by the
platform's authors with known expected outcomes, and no pre-registered
hypothesis test was run. The frozen-artifact discipline makes the
workloads auditable but not independent. Approval quality (whether human
approvers scrutinize records) is not measured; the replication episode's
scripted approver avoids claiming otherwise. Receipt grading under a
mis-applying device (a transport acknowledgment where the value is never
applied) is covered by design and by disconnect-fault checks, but no
experiment forced that exact failure; the readback-mismatch path
(re-reads, downgrade to unverified) is exercised by the
governance suite rather than by a dedicated plant-level probe.
Configuration governance is a further open side door: limits and recipe
windows are themselves configuration, and putting configuration changes
under the same approval-and-ledger discipline (management of change) is
future work. So is quantitative ramp-up guidance for the
manual-to-auto promotion path of
Section~\ref{sec:discussion-practitioners}.

\textbf{Reliability.} The frozen-baseline mechanism reproduced the
ablation bit-exactly across twelve days but correctly refused
reproduction on a drifted scenario, so regression claims are scoped to
commit- and hash-matched runs.

\textbf{Known product gaps.} The case study surfaced a real assembly
gap: the dataset builder derives control dimensions from acquisition
nodes only, so setpoint trajectories that live in the write ledger do
not yet flow into training data without an additional readback node.
The affected training path was locked closed by the twin gate rather
than producing an unsafe model. The audit ledger is append-only within
the platform process; externalizing it to write-once storage is future
work.
